\documentclass[10pt,a4paper]{amsart}
\usepackage[margin=19mm]{geometry}
\usepackage{amsmath,amssymb,mathtools}
\usepackage[scr=rsfs,cal=euler]{mathalfa}
\usepackage{xfrac}
\usepackage[unicode,hidelinks,bookmarks=false]{hyperref}
\newcommand{\ie}{i.e.\ }
\newcommand{\cf}{cf.\ }
\newcommand{\resp}{respectively\ }
\newcommand{\Subsection}{\subsection}
\providecommand{\nobreakdash}{\nobreak\hbox{-}}
\setlength{\emergencystretch}{2em}
\multlinegap0pt
\usepackage{bm}
\usepackage{bbm}
\usepackage{dsfont}
\usepackage{xspace}
\usepackage{cancel}
\usepackage{mathtools}
\usepackage{amsthm}
\makeatletter
\@ifundefined{remark}{\newtheorem{remark}{Remark}}{}
\@ifundefined{proposition}{\newtheorem{proposition}[remark]{Proposition}}{}
\@ifundefined{theorem}{\newtheorem{theorem}[remark]{Theorem}}{}
\makeatother
\title[Independent domination in central graphs]{Independent domination in central graphs}
\author{Abel Cabrera-Mart\'inez, Jos\'e Luis L\'opez-Carmona, Ismael Rios-Villamar, Alejandro Serrano-D\'iaz}
\date{Source: arXiv:2609.16357v1 (CC BY 4.0)}
\begin{document}
\maketitle

\noindent\emph{Source and licence.} Independent domination in central graphs. Version arXiv:2609.16357v1 (CC BY 4.0), distributed under the Creative Commons Attribution 4.0 International licence, as recorded in the arXiv licence metadata for this exact version and shown on the abstract page https://arxiv.org/abs/2609.16357v1. Full rights notice: licenses/arxiv-2609.16357-CC-BY-4.0.txt.\par\smallskip

\noindent\textit{Editorial scope.} Counted unit: the Theorem whose source label is \texttt{None} in the article (source0.tex lines 545--548, source label \texttt{None}), delivered together with the article's own complete original proof of it (source0.tex lines 550--626, one \texttt{proof} environment of 77 lines). No mathematics is written, repaired or re-derived: statement and proof are the article's own text line for line. Required context retained uncounted: theo-i-C(G) (lines 195--198); theo-lower-upper-bounds (lines 360--363); prop-clique (lines 535--538). Every definition, display and label of those lines is retained unchanged. Omitted from this package and not redistributed: 1--194, 199--359, 364--534, 539--544, 549--549, 627--693. Nothing inside a retained excerpt is omitted or altered: every retained body is delivered exactly as the article's lines are written, so no omission receipt of any kind accompanies this package. Citations: the retained excerpts carry no citation command at all, so no bibliography entry of the article is transcribed and no reference is redistributed. Reference adaptation: none was needed and none was made; every reference inside the delivered unit resolves inside it, so no unresolved reference or citation is produced. Printed slips kept literal: no printed word, symbol, hypothesis or number of the counted statement or of its proof was altered. Layout and class: the article's own class is \texttt{article} with packages including geometry, amssymb, amsmath, amsthm, inputenc, graphicx, hyperref, color, enumerate, enumitem, tikz, xifthen, verbatim; the wrapper is the lane's pinned \texttt{amsart} editorial wrapper at 10pt on A4, which restates the theorem-like environments the retained text needs and adds the article's own macro definitions that the retained text needs (none), with extra wrapper packages (bm, bbm, dsfont, xspace, cancel, mathtools). The delivered numbering is the unit's own. Assets: no retained span contains a figure environment, a graphics-inclusion command, a PDF-page inclusion, a table or a TikZ picture; no image file is copied into this package, the package declares no asset dependency and no image file is redistributed with it. Licence: version arXiv:2609.16357v1 of this article is distributed under the Creative Commons Attribution 4.0 International licence, as recorded in the arXiv licence metadata for this exact version and shown on the abstract page https://arxiv.org/abs/2609.16357v1. A full rights notice is delivered at licenses/arxiv-2609.16357-CC-BY-4.0.txt. Characters: every retained excerpt is pure ASCII, so the delivered engine, XeLaTeX with its default Latin Modern text font, reports no missing character.
\begin{theorem}\label{theo-i-C(G)}
Let $G$ be a connected graph of order $n\geq 3$. Then
$$i (\mathtt{C}(G)) = \min \{ |X| + |E(G-X)| + \delta(P_{G}(X)): X\in \mathcal{C}l(G)\}.$$
\end{theorem}

\begin{theorem}\label{theo-lower-upper-bounds}
Let $G$ be a connected graph of size $m\geq 2$ which is not a complete graph. Then
$$m- \frac{\omega(G)(2\Delta(G)-\omega(G)-1)}{2}\leq i(\mathtt{C}(G))\leq m- \frac{\omega(G)(2\delta(G)-\omega(G)-1)}{2}.$$
\end{theorem}

\begin{proposition}\label{prop-clique}
For any nontrivial connected graph $G$ which is not a complete graph,
$$\omega(\mathtt{C}(G)) = \alpha(G).$$
\end{proposition}

\begin{theorem}
Let $G$ be a connected graph of order $n\geq 3$ and size $m$. Then
$$i (\mathtt{C}^{2}(G)) = m + \frac{n(n-1)}{2} + \frac{\alpha(G)^{2} - \alpha(G) (2n-3)}{2}.$$
\end{theorem}

\begin{proof}
First, we observe that 
\begin{equation}\label{eq-C2-A1}
|E(\mathtt{C}(G))| = 2|E(G)| + |E(\overline{G})| = m + \frac{n(n-1)}{2}.
\end{equation}
Let $I$ be an $\alpha(G)$-set. It is straightforward to verify that $I \in \mathcal{C}l (\mathtt{C}(G))$, and by Proposition~\ref{prop-clique} we have that $|I|=\omega(\mathtt{C}(G))$. Consequently, $I$ is also an $\omega(\mathtt{C}(G))$-set.  The following equalities follow from a counting argument on the number of edges in $\mathtt{C}(G)-I$, together with the fact that $\sum_{v\in I} |N_{\mathtt{C}(G)} (v)| = \alpha(G) (n-1)$ and \eqref{eq-C2-A1}.
\begin{equation}\label{eq-C2-A2}
\begin{split}
|E(\mathtt{C}(G)-I)| &= |E(\mathtt{C}(G))| -  |E_{\mathtt{C}(G)} (I,I)| - |E_{\mathtt{C}(G)} (I, V(\mathtt{C}(G)) \setminus I)| \\[5pt]
&= m + \frac{n(n-1)}{2} - \sum_{v\in I} |N_{\mathtt{C}(G)} (v)| + \frac{\alpha(G)(\alpha(G)-1)}{2} \\[5pt]
&= m + \frac{n(n-1)}{2} + \frac{\alpha(G)^{2} -\alpha(G) (2n-1)}{2}.
\end{split}
\end{equation}
Since $I$ is an $\omega(\mathtt{C}(G))$-set, it follows that $\delta(P_{\mathtt{C}(G)}(I)) = 0$. From Theorem~\ref{theo-i-C(G)} and \eqref{eq-C2-A2} we obtain that
\begin{equation}\label{eq-C2-up}
\begin{split}
i (\mathtt{C}^{2}(G)) &\leq |I| + |E(\mathtt{C}(G)-I)| + \delta(P_{\mathtt{C}(G)}(I)) \\[5pt]
&= \alpha (G) + m + \frac{n(n-1)}{2} + \frac{\alpha(G)^{2} -\alpha(G) (2n-1)}{2}\\[5pt]
&= m + \frac{n(n-1)}{2} + \frac{\alpha(G)^{2} -\alpha(G) (2n-3)}{2}.
\end{split}
\end{equation}

\vspace{.2cm}

\noindent
Finally, we prove that \eqref{eq-C2-up} consists of a chain of equalities. To this end, we consider the following two cases.

\vspace{.2cm}

\noindent
Case 1: $G\not\cong K_n$. By Proposition \ref{prop-clique} we have that $\omega (\mathtt{C}(G)) = \alpha(G)$. By the lower bound given in Theorem~\ref{theo-lower-upper-bounds}, the fact that $\Delta(\mathtt{C}(G)) = n-1$ and \eqref{eq-C2-A1} we obtain that
\begin{equation}\label{eq-C2-low}
\begin{split}
i (\mathtt{C}^{2}(G)) &\geq |E(\mathtt{C}(G))| - \frac{\omega(\mathtt{C}(G))(2\Delta(\mathtt{C}(G))- \omega(\mathtt{C}(G)) -1 )} {2} \\[6pt]
&= m + \frac{n(n-1)}{2} + \frac{\alpha(G)^{2} - \alpha(G) (2n-3)}{2}.
\end{split}
\end{equation}
Therefore, from \eqref{eq-C2-up} and \eqref{eq-C2-low} it follows that
$$i (\mathtt{C}^{2}(G)) = m + \frac{n(n-1)}{2} + \frac{\alpha(G)^{2} - \alpha(G) (2n-3)}{2}.$$

\vspace{.2cm}

\noindent
Case 2: $G \cong K_n$. By Theorem~\ref{theo-i-C(G)}, there exists a set $S \in \mathcal{C}l (\mathtt{C}(G))$ such that 
$i(\mathtt{C}^{2}(G)) = |S| + |E(\mathtt{C}(G)-S)| + \delta(P_{\mathtt{C}(G)}(S))$. It is straightforward to verify that $|S|\leq 2$. Since $n\geq 3$, it follows that $\delta(P_{\mathtt{C}(G)}(S)) = 0$. Hence,
\begin{equation}\label{eq-C2-A3}
i(\mathtt{C}^{2}(G)) = |S| + |E(\mathtt{C}(G)-S)|.
\end{equation}
Next, we consider the following three complementary subcases.

\vspace{.2cm}
		
\noindent
Subcase 2.1: $|S| = 1$ and $S \subseteq V_E(G)$. Observe that $|E(\mathtt{C}(G)-S)| = n(n-1) - 2$. Substituting this equality into \eqref{eq-C2-A3}, we obtain that $i(\mathtt{C}^{2}(G))= n(n-1)-1$.
If $n\geq 4$, then we obtain a contradiction to \eqref{eq-C2-up}. Hence, $n=3$ and consequently, $$i(\mathtt{C}^{2}(G)) = 5 = m + \frac{n(n-1)}{2} + \frac{\alpha(G)^{2} -\alpha(G) (2n-3)}{2}.$$ 

\vspace{.2cm}
	
\noindent
Subcase 2.2: $|S| = 1$ and $S \subseteq V(G)$. Observe that $|E(\mathtt{C}(G)-S)| = (n-1)^{2}$. Substituting this equality into \eqref{eq-C2-A3}, we obtain that
\begin{equation*}
i(\mathtt{C}^{2}(G)) = 1 + (n-1)^{2}=m + \frac{n(n-1)}{2} + \frac{\alpha(G)^{2} -\alpha(G) (2n-3)}{2}.
\end{equation*}

\vspace{.2cm}

\noindent
Subcase 2.3: $|S| = 2$. Observe that $|E(\mathtt{C}(G)-S)| = (n-1)^{2} - 1$. Substituting this equality into \eqref{eq-C2-A3}, we obtain that
\begin{equation*}
i(\mathtt{C}^{2}(G)) = 1 + (n-1)^{2}=   m + \frac{n(n-1)}{2} + \frac{\alpha(G)^{2} -\alpha(G) (2n-3)}{2}
\end{equation*}

\vspace{.25cm}

\noindent
From the previous two cases, we conclude that $i(\mathtt{C}^{2}(G)) = m + \frac{n(n-1)}{2} + \frac{\alpha(G)^{2} -\alpha(G) (2n-3)}{2}$, which completes the proof.
\end{proof}

\end{document}
